\scanpage{188}
\mspage{62}
% archive sheet number: 187

Je vais quand même esquisser la démonstration. Soit $X=\varinjlim_I X_i$, $I$ filtrant, les $X_i$ dans $M$, il faut prouver que $X\in M$ (moyennant la \inserted{l'hyp. de} stabilité sur $M^{\unsure{\pi|\Pi}}=M\cap\mathrm{Filt}^{\Pi}$).

\diagram{\begin{tikzcd}[ampersand replacement=\&,column sep=small,row sep=small]
 \& C\arrow[r]\arrow[d]\arrow[dl,dashed,"{\unsure{z|3|\mathfrak z}}"'] \& I\arrow[d,"{\unsure{\varphi|\phi}}"]\\
 \underline M^\pi \& \underline M^\pi/\underline M\arrow[l,dashed,"{\text{``source''}}"]\arrow[r,"{\text{``but''}}"'] \& \underline M
 \arrow[ul,phantom,"{\text{cart}}" description]
\end{tikzcd}}

$C$ est une catégorie cofibrée sur $I$, dont les fibres $C_i$ sont

\[
 C_i\simeq\underline M^\pi/X_i,
\]

or comme $X_{\unsure{i|0}}\in M$, on voit, à cause de la condition de $\pi$-accessibilité, que $M^\pi/X_i$ est \uline{filtrante}. \deleted{\unsure{Cela} En effet} \inserted{\unsure{Car ce fait|Car ceci}} \deleted{\unsure{Lemme} \unsure{de ce fait}} \inserted{dégageons un}

\margin{\rule{0.4pt}{1.4cm}}\uline{Lemme} 5.\deleted{\unsure{9}}10. Soit $M$ une partie $\pi$-accessible de $\mathcal M$. Alors pour tt $\deleted{\unsure{Z}}Y\in\deleted{\unsure{X}}M$, $M^\pi/\deleted{\unsure{Z}}Y$ est une catégorie \uline{filtrante}, et $\mathrm{Ss}_\pi(Y)$ y est cofinale.

\deleted{En effet, les $Z\longrightarrow Y$ avec $Z$ dans $M^\pi$ (ou seulement dans $\mathrm{Filt}^{\pi}$) se factorisent par $Z\longrightarrow Z_\alpha\hookrightarrow Y$, avec $Z_\alpha$ dans $\mathrm{Ss}_\pi(Y)$, or \deleted{\unsure{les}} si $Y\in M$, les $Z_\alpha$ tels que $Z_\alpha\in M$ sont cofinaux, donc les $Z\longrightarrow Y$ ($Z\in M^\pi$) se factorisent par $Z_\alpha\lhook\joinrel\longrightarrow Y$ avec $Z_\alpha\in M^\pi$. Cela montre que si $Z\to Y,Z'\to Y$ sont dans $M^\pi/Y$, deux obj sont majorés par un $Z_\alpha\to Y$ $\in\mathrm{Ss}_\pi(Y)\unsure{/Y}\subset M^\pi/Y$. D'autre}

